\bigskip\noindent{\Large\bfseries\color{apren} Test 2. Forma polar, potències i arrels}\par\medskip
\iniciTest{complexosT2}{a,b,a,c,b,a,c,a,b,c}
\begin{enumerate}
\pregunta Si $z=r_\alpha$ i $w=s_\beta$, quin és el producte $zw$ en forma polar?
\begin{enumerate}
\opcio{a}{$rs_{\alpha+\beta}$}
\opcio{b}{$(r+s)_{\alpha+\beta}$}
\opcio{c}{$rs_{\alpha-\beta}$}
\opcio{d}{$\frac r s{}_{\alpha+\beta}$}
\end{enumerate}\feedback
\pregunta Si $z=r_\alpha$ i $w=s_\beta$ amb $w\ne0$, quin és $z/w$?
\begin{enumerate}
\opcio{a}{$\frac r s{}_{\alpha+\beta}$}
\opcio{b}{$\frac r s{}_{\alpha-\beta}$}
\opcio{c}{$(r-s)_{\alpha-\beta}$}
\opcio{d}{$rs_{\alpha/\beta}$}
\end{enumerate}\feedback
\pregunta La fórmula de De Moivre afirma que:
\begin{enumerate}
\opcio{a}{$[r(\cos\alpha+i\sin\alpha)]^n=r^n(\cos n\alpha+i\sin n\alpha)$}
\opcio{b}{$z^n=nz$}
\opcio{c}{$|z^n|=n|z|$}
\opcio{d}{$\arg(z^n)=\arg z/n$}
\end{enumerate}\feedback
\pregunta Quantes arrels cinquenes té un complex no nul?
\begin{enumerate}
\opcio{a}{Una}
\opcio{b}{Dues}
\opcio{c}{Cinc}
\opcio{d}{Infinites}
\end{enumerate}\feedback
\pregunta Els arguments de dues arrels $n$-èsimes consecutives d’un complex difereixen en:
\begin{enumerate}
\opcio{a}{$\pi/n$}
\opcio{b}{$2\pi/n$}
\opcio{c}{$n\pi$}
\opcio{d}{$360n^\circ$}
\end{enumerate}\feedback
\pregunta Quina és la forma trigonomètrica de $r_\alpha$?
\begin{enumerate}
\opcio{a}{$r(\cos\alpha+i\sin\alpha)$}
\opcio{b}{$r+i\alpha$}
\opcio{c}{$\alpha(\cos r+i\sin r)$}
\opcio{d}{$r\cos\alpha+i\sin r$}
\end{enumerate}\feedback
\pregunta En forma polar, la suma de dos complexos es calcula:
\begin{enumerate}
\opcio{a}{sumant mòduls i arguments}
\opcio{b}{restant arguments}
\opcio{c}{passant normalment a forma binòmica}
\opcio{d}{multiplicant mòduls}
\end{enumerate}\feedback
\pregunta Si $z=2_{30^\circ}$, aleshores $z^3$ és:
\begin{enumerate}
\opcio{a}{$8_{90^\circ}$}
\opcio{b}{$6_{90^\circ}$}
\opcio{c}{$8_{10^\circ}$}
\opcio{d}{$2_{90^\circ}$}
\end{enumerate}\feedback
\pregunta Les arrels $n$-èsimes d’un complex no nul tenen afixos situats:
\begin{enumerate}
\opcio{a}{sobre una recta}
\opcio{b}{en els vèrtexs d’un polígon regular}
\opcio{c}{sempre sobre l’eix real}
\opcio{d}{sempre en el primer quadrant}
\end{enumerate}\feedback
\pregunta Si $z=r_\alpha$, quin és el mòdul de $z^{-1}$?
\begin{enumerate}
\opcio{a}{$r$}
\opcio{b}{$-r$}
\opcio{c}{$1/r$}
\opcio{d}{$r^2$}
\end{enumerate}\feedback
\end{enumerate}
